\title{The Era of 1-bit LLMs: All Large Language Models are in 1.58 Bits}

\begin{abstract}
Recent advancements like BitNet~\cite{bitnet} are ushering in a transformative phase for 1-bit Large Language Models (LLMs). Here, we propose a new 1-bit LLM variant, \textbf{\text{BitNet b1.58}}, where each individual parameter (or weight) of the LLM adopts a ternary representation from the set \{-1, 0, 1\}. This model achieves parity with full-precision (FP16 or BF16) Transformer LLMs of comparable size and trained with an equivalent number of tokens, demonstrating competitive perplexity and performance on downstream tasks. In addition, it delivers pronounced efficiency gains in latency, memory usage, throughput, and energy demand. More notably, the introduction of the 1.58-bit LLM establishes a novel scaling law and methodological framework for training future generations of LLMs that harmonize high performance with efficiency. It also catalyzes new computational strategies, laying the groundwork for dedicated hardware tailored for 1-bit LLMs.
\end{abstract}

\section{The Era of 1-bit LLMs}

The landscape of artificial intelligence has recently experienced substantial advancement in both the magnitude and functionality of Large Language Models (LLMs). Although these models excel at various natural language processing tasks, their escalating scale introduces considerable obstacles to practical deployment, with significant concerns emerging about their energy requirements and associated environmental and economic burdens.

To mitigate these difficulties, researchers have turned to post-training quantization as a strategy for generating low-bitwidth models for inference~\cite{smoothquant, gptq, quip, quip_sharp}. This method lessens the numerical precision of weights and activations, resulting in considerable savings in memory footprint and computational load for LLMs. The progression has moved beyond 16-bit down to models with only 4 bits~\cite{gptq, awq}. Nonetheless, post-training quantization—despite its broad adoption in commercial LLMs—remains a sub-optimal solution.

Recent advances in 1-bit model designs, exemplified by BitNet~\cite{bitnet}, have shown substantial potential for lowering the operational costs of LLMs without significant loss in performance. Conventional LLMs utilize 16-bit floating-point representations (FP16 or BF16) and depend heavily on matrix multiplication, which is computationally intensive due to repeated floating-point additions and multiplications. BitNet, in contrast, performs its matrix multiplications using only integer addition, dramatically reducing the energy consumption associated with running LLMs. Given that power efficiency often constrains compute performance on hardware accelerators, the reduced energy requirements can also accelerate computational throughput.

Another significant bottleneck arises from the need to transfer large volumes of model parameters from DRAM to on-chip memory (e.g., SRAM) during inference, which can be resource-intensive. Expanding SRAM capacities has been proposed as a solution to boost throughput, but this incurs higher costs than DRAM. By comparison, 1-bit LLMs occupy less memory overall—both in storage capacity and bandwidth—than their full-precision counterparts. This compactness leads to quicker and less costly parameter loading, thereby increasing the efficiency and speed of inference.

Within this context, we present \textbf{\text{BitNet b1.58}}, an important advancement among 1-bit LLMs. In this variant, each parameter can be one of three values: \{-1, 0, 1\}, incorporating a zero value into the traditional 1-bit BitNet design and thereby achieving 1.58 bits per parameter. \text{BitNet b1.58} preserves the key strengths of the original 1-bit BitNet, such as its highly optimized computational paradigm that nearly eliminates the need for multiplication in matrix operations, enabling significant energy and memory savings when compared to FP16-based LLMs. On top of these efficiencies, \text{BitNet b1.58} introduces further enhancements. First, the model’s expressivity improves through explicit feature filtering enabled by zero-valued weights, boosting performance over traditional 1-bit LLMs. Second, empirical results reveal that \text{BitNet b1.58} can attain parity with full-precision (FP16) models in terms of perplexity and downstream task effectiveness, beginning at the 3B parameter scale, provided the architectural and training settings are kept consistent.

\section{BitNet b1.58}

\text{BitNet b1.58} is constructed upon the BitNet framework, where the standard \emph{nn.Linear} layers of Transformers are substituted with \emph{BitLinear} layers. The model is initialized from the beginning, featuring weights quantized to 1.58 bits and activations represented with 8 bits. Several changes have been incorporated relative to the initial BitNet architecture, as outlined below.

\paragraph{Quantization Function.} To limit weights to the set $\{-1, 0, +1\}$, we employ an \emph{absmean} quantization technique. This process first normalizes the weight matrix by dividing by its mean absolute value and then quantizes each element to the closest integer within $\{-1, 0, +1\}$:

\begin{equation}
    \widetilde{W} = \mathrm{RoundClip}(\frac{W}{\gamma + \epsilon}, -1, 1),
\end{equation}

\begin{equation}
    \mathrm{RoundClip}(x, a, b) = \max(a, \min(b, \mathrm{round}(x))),
\end{equation}

\begin{equation}
    \gamma = \frac{1}{nm}\sum_{ij} |W_{ij}|.
\end{equation}

For activations, the quantization approach mirrors that of BitNet, except that activations are not rescaled prior to the non-linearities into $[0, Q_b]$. Instead, scaling occurs to map all activations per token to the interval $[-Q_b, Q_b]$, which eliminates zero-point quantization. This adaptation not only simplifies implementation and system-level optimizations, but also, in our empirical studies, has only an insignificant impact on performance.

\paragraph{LLaMA-alike Components.}

Given that LLaMA~\cite{llama, llama2} has emerged as a standard reference model within open-source LLMs, \text{BitNet b1.58} is constructed with architectural elements consistent with LLaMA to support integration within the open-source ecosystem. Concretely, it incorporates RMSNorm~\cite{rmsnorm}, SwiGLU~\cite{swiglu}, and rotary embedding~\cite{rope}, and eliminates all bias terms. This approach facilitates straightforward adoption of \text{BitNet b1.58} within widely used open-source tools, such as Huggingface, vLLM~\cite{vllm}, and llama.cpp\footnote{\url{https://github.com/ggerganov/llama.cpp}}, thus minimizing integration overhead.

\section{Results}

We evaluated \text{BitNet b1.58} against our implementation of FP16 \text{LLaMA LLM} across several model scales. Both models underwent pre-training on the RedPajama dataset~\cite{redpajama}, utilizing 100 billion tokens to maintain an equitable assessment. To gauge their capabilities, we examined zero-shot outcomes on an assortment of language benchmarks: ARC-Easy~\cite{arc}, ARC-Challenge~\cite{arc}, Hellaswag~\cite{hellaswag}, Winogrande~\cite{winoGrande}, PIQA~\cite{piqa}, OpenbookQA~\cite{openbookqa}, and BoolQ~\cite{boolq}. Additionally, validation perplexities were recorded on WikiText2~\cite{wikitext2} and C4~\cite{c4} datasets.

We also assessed both the GPU memory consumption and inference latency of \text{LLaMA LLM} and \text{BitNet b1.58}. Our measurements leveraged the FasterTransformer\footnote{\url{https://github.com/NVIDIA/FasterTransformer}} framework, recognized for its optimized LLM inference capabilities on GPU platforms. For \text{BitNet b1.58}, the 2-bit kernel provided by Ladder~\cite{ladder} was employed. The principal metric reported was time per output token, as this constitutes the main computational bottleneck during inference.

In Table~\ref{tab:ppl}, we present a summary of perplexity and resource costs for both \text{BitNet b1.58} and \text{LLaMA LLM}. The data indicate that, from the 3B model scale onwards, \text{BitNet b1.58} achieves comparable perplexity to the FP16 \text{LLaMA LLM}, while attaining 2.71$\times$ faster inference and reducing GPU memory usage by a factor of 3.55. Notably, with 3.9B parameters, \text{BitNet b1.58} demonstrates a 2.4$\times$ speedup, a 3.32$\times$ decrease in memory demands, and delivers superior performance relative to the 3B version of \text{LLaMA LLM}.

The zero-shot evaluation results, detailed in Table~\ref{tab:zero-shot}, were obtained using the evaluation framework provided by \emph{lm-evaluation-harness}\footnote{\url{https://github.com/EleutherAI/lm-evaluation-harness}}. These findings illustrate that as model size increases, the difference in performance between \text{BitNet b1.58} and \text{LLaMA LLM} diminishes. Crucially, from the 3B parameter scale, \text{BitNet b1.58} equals the accuracy of the full-precision baseline. Corroborating the perplexity trends, the end-task accuracies reveal that the 3.9B \text{BitNet b1.58} variant not only consumes less memory and achieves lower latency but also surpasses the \text{LLaMA LLM} 3B model. These results establish that \text{BitNet b1.58} represents a Pareto improvement compared to leading large language models.

\paragraph{Memory and Latency}

To investigate the scalability of our approach, we increased the model size to 7B, 13B, and 70B and analyzed the associated computational cost. As depicted in Figure~\ref{fig:latency-memory-energy}, both latency and memory usage exhibit patterns indicating enhanced speed-up with larger model sizes. Notably, \text{BitNet b1.58} at the 70B parameter scale achieves a speed 4.1 times greater than the \text{LLaMA LLM} reference model. This improvement arises because the computational demand for \emph{nn.Linear} operations escalates with model size. The observed pattern in memory usage can be attributed to the embeddings, which retain full precision; their share of total memory decreases as models grow larger. Latency and memory metrics were obtained using a 2-bit kernel, suggesting that further cost reductions are possible through additional optimizations.

\paragraph{Energy}

We also performed an estimation of the energy expenditure from arithmetic operations in both \text{BitNet b1.58} and \text{LLaMA LLM}. Emphasis was placed on the matrix multiplication component, given its significant impact on overall LLM cost. As illustrated in Figure~\ref{fig:energy}, \text{BitNet b1.58} primarily relies on INT8 addition operations, whereas \text{LLaMA LLM} incorporates both FP16 addition and multiplication. Following the energy modeling in~\cite{energycost, pokebnn}, our analysis indicates that \text{BitNet b1.58} achieves a 71.4-fold reduction in arithmetic operation energy for matrix multiplication on 7nm hardware. Additionally, we measured the total energy consumption for models processing sequences of 512 tokens. The findings reveal that as model scale increases, \text{BitNet b1.58} demonstrates growing energy efficiency advantages over the FP16 \text{LLaMA LLM}. This trend is attributable to the increased proportion of computation devoted to \emph{nn.Linear} operations with larger models, while the relative contribution from other components becomes less significant.

\paragraph{Throughput}

We evaluate and contrast the throughput of \text{BitNet b1.58} and \text{LLaMA LLM}—both with 70B parameters—deployed across two 80GB A100 GPUs, utilizing pipeline parallelism~\cite{gpipe} to facilitate the execution of \text{LLaMA LLM} 70B on the hardware. To identify maximum capacity, we incrementally increased the batch size until the GPU memory was saturated, while keeping the sequence length fixed at 512. The data summarized in Table~\ref{tab:throughput} reveals that \text{BitNet b1.58} 70B accommodates up to 11 times the batch size achievable by \text{LLaMA LLM}, translating to an 8.9-fold gain in throughput.

\textbf{\text{BitNet b1.58} establishes a new scaling paradigm relating to the interplay of model size, performance, and inference efficiency}. Leveraging the outcomes illustrated in Figure~\ref{fig:latency-memory-energy} and Figure~\ref{fig:energy}, we observe the following correspondences between 1.58-bit models and their 16-bit counterparts:

\begin{itemize}
    \item A 13B parameter BitNet b1.58 outperforms a 3B FP16 LLM concerning latency, memory requirements, and energy usage.
    \item A 30B BitNet b1.58 demonstrates superior efficiency in latency, memory footprint, and energy consumption relative to a 7B FP16 LLM.
    \item A 70B BitNet b1.58 shows improved efficiency in all the aforementioned aspects when compared with a 13B FP16 LLM.
\end{itemize}

\paragraph{Training with 2T Tokens}

The token count used for training is a key determinant of LLM quality. To assess how well \text{BitNet b1.58} scales with token count, we trained a \text{BitNet b1.58} model on 2T tokens, employing the same data protocol as StableLM-3B~\cite{StableLM-3B-4E1T}, a leading open-source 3B parameter model. Both models underwent evaluation using a benchmark suite composed of Winogrande~\cite{winoGrande}, PIQA~\cite{piqa}, SciQ~\cite{sciq}, LAMBADA~\cite{lambada}, and ARC-easy~\cite{arc}. Zero-shot accuracies are reported in Table~\ref{tab:2t}. Where both accuracy and normalized accuracy apply, we averaged these metrics. For reference, StableLM 3b results at 2T tokens are sourced from its official technical documentation. The evidence indicates that \text{BitNet b1.58} consistently surpasses StableLM 3b across all evaluated tasks, suggesting that LLMs quantized to 1.58 bits also maintain robust generalization ability.

\section{Discussion and Future Work}

\textbf{1-bit Mixture-of-Experts (MoE) LLMs}

The Mixture-of-Experts (MoE) architecture has established itself as an efficient and scalable solution for large language models (LLMs), primarily due to its reduction of computational FLOPs. However, MoE models are often hampered by substantial memory requirements and significant inter-chip communication costs, which impede their practicality for large-scale deployments. 1.58-bit LLMs offer promising avenues for resolving these concerns. By reducing memory usage, fewer hardware units are necessary to host MoE models, thereby cutting down the complexity and cost of deployment. This optimization also greatly lessens the communication overhead involved in transmitting activations among devices. Ideally, if model size permits placement on a single chip, these overheads could be entirely eliminated.

\textbf{Native Support of Long Sequence in LLMs}

As the significance of LLMs grows, efficiently processing lengthy input sequences has become a fundamental challenge. One of the primary obstacles in long sequence inference lies in the considerable memory demand for key-value (KV) caches. \text{BitNet b1.58} makes progress in addressing this issue by compressing activations from the traditional 16 bits to just 8 bits, effectively doubling the feasible context length for the same hardware resources. There is also scope for further lossless compression to 4 bits or even lower with 1.58-bit LLMs, a direction intended for continued research.

\textbf{LLMs on Edge and Mobile}

Deploying LLMs with 1.58-bit quantization opens new horizons for running advanced language models on edge and mobile platforms. These devices generally contend with stringent memory and compute restrictions, constraining both the size and performance of LLMs. By minimizing both memory usage and power consumption, 1.58-bit LLMs enable the deployment of robust models on hardware that was previously insufficient. This enhancement can expand the range and utility of edge and mobile AI applications. Furthermore, 1.58-bit LLMs align well with the capabilities of CPUs, which are prevalent in mobile and edge systems, thereby making \text{BitNet b1.58} suitable for efficient and performant inference in these contexts.

\textbf{New Hardware for 1-bit LLMs}

Recent advancements, exemplified by systems such as Groq\footnote{\url{https://groq.com/}}, highlight the advantages of dedicated hardware—including language processing units (LPUs)—optimized for LLM workloads. Building on this momentum, there is a strong case for the development of next-generation hardware and computational architectures tailored to the distinct demands of 1-bit LLMs, as enabled by the BitNet approach~\cite{bitnet}.

\end{document}